\begin{definition}[BNT multiplicity-spectrum comparison]%
    \label{def:mpdo_bnt_multiplicity_spectrum_comparison}
    \lean{MPOTensor.BNTMultiplicitySpectrumComparison}
    \leanok
    \uses{def:mpdo_diagonal_chi_family, def:mpdo_bnt_label_trace_scalar_family}
    Let $\mu_\alpha$ and $\nu_\gamma$ be diagonal matrices with complex
    entries, representing one-site and two-site multiplicities, respectively.
    A multiplicity-spectrum comparison requires
    $\tr(\mu_\alpha)=m_\alpha$ and, for every $\gamma$, equality with
    multiplicity between the entries of $\nu_\gamma$ and the products
    $\mu_{\alpha,i}\mu_{\beta,j}\chi_{\alpha,\beta,\gamma,k}$.
    For the positive diagonal matrices in the vertical canonical
    decompositions, this is the conclusion of the BNT comparison and positive
    power-sum argument in
    \cite[Appendix~C.4, lines~2048--2058]{Cirac2016MPDO_arXiv}.
    The definition itself does not require positivity or a realization by
    an MPDO tensor.
\end{definition}

\begin{definition}[Two-site multiplicity spectrum]%
    \label{def:mpdo_bnt_two_site_multiplicity_spectrum}
    \lean{MPOTensor.BNTAlgebraTensorClause.TwoSiteMultiplicitySpectrum}
    \leanok
    \uses{def:mpdo_bnt_algebra_tensor_clause}
    Let $M$ have a vertical decomposition
    $\widetilde M=\bigoplus_\alpha\mu_\alpha\otimes M_\alpha$
    with a basis of normal tensors $\{M_\alpha\}_\alpha$ and positive diagonal matrices
    $\mu_\alpha$, satisfying the BNT algebra condition.
    A two-site multiplicity spectrum consists of a vertical canonical
    decomposition
    \begin{align}
        \widetilde{M^{(2)}}
         & =\bigoplus_\delta \nu_\delta\otimes A_\delta,
        \notag
    \end{align}
    a bijection $\sigma$ from the one-site sectors to the two-site sectors such
    that $M_\gamma$ and $A_{\sigma(\gamma)}$ generate the same matrix product
    vectors at every positive length, and, for every $\gamma$, the multiset
    equality
    \begin{align}
        \bigl\{\nu_{\sigma(\gamma),r}:r\bigr\}
         & =
        \bigl\{\mu_{\alpha,i}\mu_{\beta,j}
        \chi_{\alpha,\beta,\gamma,k}:\alpha,\beta,i,j,k\bigr\}.
        \notag
    \end{align}
    This is the decomposition and comparison used in
    \cite[Appendix~C.4, lines~2046--2058]{Cirac2016MPDO_arXiv}.
\end{definition}

\begin{definition}[Exact invertible gauges for the two-site sectors]%
    \label{def:mpdo_bnt_two_site_exact_sector_gauge}
    \lean{MPOTensor.BNTAlgebraTensorClause.TwoSiteExactSectorGauge}
    \leanok
    \uses{def:mpdo_bnt_two_site_multiplicity_spectrum}
    Exact invertible gauges for a two-site multiplicity spectrum require,
    for every one-site sector $\gamma$, equality
    of the paired bond dimensions $d_\gamma=d_{\sigma(\gamma)}$ and an
    invertible matrix $Z_\gamma\in\GL(d_\gamma,\C)$ such that, for every
    physical index $i$ of the vertical tensor,
    \begin{align}
        A_{\sigma(\gamma)}^i
         & =Z_\gamma M_\gamma^i Z_\gamma^{-1}.
        \notag
    \end{align}
    Thus the matched normal tensors have no residual scalar phase. This is the
    exact invertible conjugacy in
    \cite[Appendix~C.4, lines~2053--2057]{Cirac2016MPDO_arXiv}.
    The definition does not require $Z_\gamma$ to be unitary.
\end{definition}

\begin{theorem}[Two-site multiplicity spectrum from the BNT algebra
        condition]%
    \label{thm:mpdo_bnt_algebra_tensor_clause_spectrum}
    \lean{MPOTensor.BNTAlgebraTensorClause.%
        toTwoSiteExactSectorGauge}
    \lean{MPOTensor.BNTAlgebraTensorClause.%
        toTwoSiteMultiplicitySpectrum}
    \leanok
    \uses{def:cpsv_canonical_form, def:mpdo, def:mpdo_bnt_algebra_tensor_clause, def:mpdo_bnt_two_site_multiplicity_spectrum, def:mpdo_bnt_two_site_exact_sector_gauge, def:mpdo_bnt_multiplicity_spectrum_comparison, def:mpdo_bnt_eventual_power_sums_multiplicity_comparison}
    Let $M$ generate MPDOs, and suppose that its MPS tensor obtained by
    pairing the physical indices is in canonical form. Suppose also that $M$ has a
    vertical decomposition
    $\widetilde M=\bigoplus_\alpha \mu_\alpha\otimes M_\alpha$
    satisfying the BNT algebra condition.
    There is a vertical canonical decomposition of the two-site blocking,
    \begin{align}
        \widetilde{M^{(2)}}
         & =\bigoplus_\delta \nu_\delta\otimes A_\delta.
        \notag
    \end{align}
    Its normal-tensor sectors are in bijection, via $\sigma$, with the sectors
    of the one-site product expansion, in which every sector occurs.
    The tensors $M_\gamma$ and $A_{\sigma(\gamma)}$ generate the same matrix
    product vectors at every positive length and, for every $\gamma$,
    \begin{align}
        \bigl\{\nu_{\sigma(\gamma),r}:r\bigr\}
         & =
        \bigl\{\mu_{\alpha,i}\mu_{\beta,j}
        \chi_{\alpha,\beta,\gamma,k}:\alpha,\beta,i,j,k\bigr\}
        \label{eq:rfp_multiplicity_spectrum}
    \end{align}
    as multisets. This is the multiplicity-spectrum comparison of
    Appendix~C.4, lines~2046--2058
    of~\cite{Cirac2016MPDO_arXiv}. The paired bond dimensions agree,
    and there are invertible matrices $Z_\gamma$ satisfying, for every
    physical index $i$ of the vertical tensor,
    \begin{align}
        A_{\sigma(\gamma)}^i
         & =Z_\gamma M_\gamma^i Z_\gamma^{-1}.
        \label{eq:rfp_sector_conjugacy}
    \end{align}
    This statement asserts exact invertible conjugacy, without requiring
    the matrices $Z_\gamma$ to be unitary.
\end{theorem}

\begin{proof}\leanok
    \uses{thm:mpdo_literal_cpsv_vertical_decomposition_block_two}
    Multiplying the two one-site decompositions and using the BNT algebra law
    expresses the two-site closed chain in terms of $O_L(M_\gamma)$, with
    coefficient
    \begin{align}
        \sum_{\alpha,\beta,i,j,k}
        \bigl(\mu_{\alpha,i}\mu_{\beta,j}
        \chi_{\alpha,\beta,\gamma,k}\bigr)^L.
        \notag
    \end{align}
    All these coefficients are positive. Hence every basis sector occurs,
    and the characterization of bases of normal tensors matches these sectors
    bijectively with those of the two-site vertical canonical decomposition.
    Normalization makes each matching scalar have modulus one, while
    positivity of two consecutive coefficients makes it positive; it is
    therefore one. The normal-tensor matching theorem also supplies equality
    of the paired bond dimensions and an invertible gauge; after the scalar is
    one, its gauge-phase relation is precisely the exact conjugacy
    in~(\ref{eq:rfp_sector_conjugacy}).
    Eventual linear independence of the basis vectors gives the corresponding
    equality of power sums. The finite positive power-sum identity then gives
    the multiset equality~(\ref{eq:rfp_multiplicity_spectrum}).
\end{proof}

\begin{theorem}[Physical compression onto a conjugate two-site sector]
    \label{thm:mpdo_bnt_exact_sector_gauge_blocked_corner}
    \lean{MPOTensor.BNTAlgebraTensorClause.TwoSiteExactSectorGauge.%
        exists_blockTwo_gauge_sector_corner}
    \lean{MPOTensor.BNTAlgebraTensorClause.TwoSiteExactSectorGauge.%
        markedChainCoefficient_gauge_eq_reflectedAdjoint}
    \leanok
    \uses{def:mpdo, def:mpdo_bnt_two_site_exact_sector_gauge}
    Suppose that the one-site and two-site sectors are related by the exact
    invertible gauges above. Let $\gamma$ be a one-site sector matched with a
    two-site sector by the bijection $\sigma$, and identify their bond
    dimensions. There are an
    isometry $V_\gamma$ from the sector bond space into the blocked physical
    space and a scalar $c_\gamma>0$ such that, for every vertical letter $v$,
    \begin{align}
        V_\gamma^\dagger\widetilde{M^{(2)}}^vV_\gamma
         & =c_\gamma Z_\gamma M_\gamma^vZ_\gamma^{-1}.
        \label{eq:mpdo_bnt_exact_sector_gauge_blocked_corner}
    \end{align}
    If, in addition, $M$ generates matrix product density operators, then for
    every pair of open bond indices and every pair of two-site tail words, the
    marked-chain coefficient of
    \begin{align}
        (Z_\gamma^\dagger Z_\gamma)M_\gamma^v
        (Z_\gamma^\dagger Z_\gamma)^{-1}
        \notag
    \end{align}
    equals the corresponding reflected-adjoint marked-chain coefficient in
    the adjoint two-site blocking. This is one of the two marked comparisons
    in Figures~7--8 of
    \cite[Proposition~4.13, lines~1909--1919]{Cirac2016MPDO_arXiv}, applied
    to the sector matching of Appendix~C.4, lines~2048--2057.

    The compression in~(\ref{eq:mpdo_bnt_exact_sector_gauge_blocked_corner})
    selects a conjugate of $M_\gamma$, not $M_\gamma$ itself. It does not
    provide an isometric compression onto a positive multiple of $M_\gamma$
    or establish the Gram identity
    \begin{align}
        (Z_\gamma^\dagger Z_\gamma)M_\gamma^v
        (Z_\gamma^\dagger Z_\gamma)^{-1}
         & =M_\gamma^v.
        \notag
    \end{align}
    This distinction is documented in
    \cite{gap:cpsv16_two_site_sector_unitary_gauge_gap}.
    The comparison with a common unblocked tail below proves the Gram
    identity without requiring such a compression of the two-site blocking.
\end{theorem}

\begin{proof}\leanok
    \uses{thm:mpdo_blocked_reference_corners, thm:mpdo_gram_dressed_reflected_marked_chain}
    Choose the first positive-weight copy of the matched two-site sector.
    Its inclusion is isometric and its compression selects the
    corresponding weighted normal tensor. Substituting the exact sector
    conjugacy gives
    (\ref{eq:mpdo_bnt_exact_sector_gauge_blocked_corner}), with the selected
    weight as $c_\gamma$. The reflected marked-chain identity for this
    compression, conjugated on its open indices by $Z_\gamma^\dagger$ and
    $(Z_\gamma^{-1})^\dagger$, gives the stated equality of marked-chain
    coefficients; the gauge factors cancel on the reflected side.
\end{proof}

\begin{theorem}[Physical compression onto a matched one-site sector]
    \label{thm:mpdo_bnt_exact_sector_one_site_raw_corner}
    \lean{MPOTensor.BNTAlgebraTensorClause.TwoSiteExactSectorGauge.%
        exists_oneSite_raw_sector_corner}
    \leanok
    \uses{def:mpdo_bnt_algebra_tensor_clause, def:mpdo_bnt_two_site_exact_sector_gauge}
    In the setting above, each sector $\gamma$ admits an isometry
    $W_\gamma\colon\C^{d_\gamma}\to\C^d$ and a scalar
    $a_\gamma>0$ such that, for every vertical letter $v$,
    \begin{align}
        W_\gamma^\dagger\widetilde M^vW_\gamma
         & =a_\gamma M_\gamma^v.
        \label{eq:mpdo_bnt_exact_sector_one_site_raw_corner}
    \end{align}
    Here the bond space of $M_\gamma$ is
    identified with that of its matched two-site sector. This compression
    selects a positive-weight copy in the one-site vertical decomposition of
    \cite[Proposition~4.13, statement lines~943--951, and Appendix~C.4,
        lines~2046--2057]{Cirac2016MPDO_arXiv}.
\end{theorem}

\begin{proof}\leanok
    \uses{thm:mpdo_retained_product_spectral_family_exists}
    Select the first copy of $M_\gamma$ in the one-site vertical direct sum.
    Its inclusion is isometric. Compressing the reconstruction identity by
    this inclusion gives
    (\ref{eq:mpdo_bnt_exact_sector_one_site_raw_corner}), with
    $a_\gamma$ equal to the positive multiplicity weight of the selected
    copy. Identify the paired bond spaces using equality of their dimensions.
\end{proof}

\begin{theorem}[Reflected marked chains with a two-site initial block]
    \label{thm:mpdo_two_site_prefix_gram_reflected_marked_chain}
    \lean{MPOTensor.IsMPDO.%
        markedChainCoefficient_blockTwoPrefix_gramDressing_eq_reflectedAdjoint}
    \leanok
    \uses{def:mpdo, def:vertical_tensor_mpdo, def:mpdo_two_site_blocking}
    Let $M$ generate matrix product density operators, let $A$ be a vertical
    tensor, and suppose that an isometry $V$ selects from the two-site
    blocking a positive multiple of an invertible conjugate of $A$.
    For some scalar $b>0$, this means that, for every vertical letter $v$,
    \begin{align}
        V^\dagger\widetilde{M^{(2)}}^vV
         & =b\,X A^vX^{-1}.
        \label{eq:mpdo_two_site_prefix_gauge_corner}
    \end{align}
    For every pair of open bond indices and every pair of equal-length words
    in the original physical alphabet, the marked-chain coefficient of
    \begin{align}
        (X^\dagger X)A^v(X^\dagger X)^{-1}
        \notag
    \end{align}
    followed by the corresponding unblocked tail of $M$ equals the
    reflected-adjoint marked-chain coefficient of $A$ followed by the
    reversed tail of $M^\dagger$. Thus the first two sites are compressed
    jointly, while all later sites remain unblocked.

    Keeping the tail unblocked supplies the common comparison needed when
    \cite[Proposition~4.13, lines~1909--1919]{Cirac2016MPDO_arXiv} is invoked
    in~\cite[Appendix~C.4, lines~2048--2057]{Cirac2016MPDO_arXiv}.
    This correction to the direct two-site comparison is recorded in
    \cite{gap:cpsv16_two_site_sector_unitary_gauge_gap}.
\end{theorem}

\begin{proof}\leanok
    \uses{thm:mpdo_reflected_marked_chain}
    Compress the first two physical sites of the density operator of length
    $N+2$ by $V$. Hermiticity exchanges the open bond indices and the two
    physical words. Taking adjoints reverses only the remaining $N$ one-site
    factors. Finally conjugate the open indices by $X^\dagger$ and
    $(X^{-1})^\dagger$. The left side becomes conjugation by $X^\dagger X$,
    while the gauge factors cancel on the reflected side.
\end{proof}

\begin{theorem}[Two-sided compression onto a matched two-site sector]
    \label{thm:mpdo_bnt_exact_sector_identity_oblique_physical_span}
    \lean{MPOTensor.twoSidedCompressionCoefficients}
    \lean{MPOTensor.linearMarkedTensor_twoSidedCompressionCoefficients}
    \lean{MPOTensor.BNTAlgebraTensorClause.TwoSiteExactSectorGauge.%
        exists_blockTwo_identity_oblique_compression}
    \lean{MPOTensor.BNTAlgebraTensorClause.TwoSiteExactSectorGauge.%
        exists_identity_physical_letter_coefficients}
    \leanok
    \uses{thm:mpdo_bnt_exact_sector_gauge_blocked_corner}
    For a matched two-site sector, choose $V_\gamma$ and
    $c_\gamma>0$ as in
    (\ref{eq:mpdo_bnt_exact_sector_gauge_blocked_corner}), and set
    \begin{align}
        L_\gamma
         & =V_\gamma(Z_\gamma^{-1})^\dagger,
        \label{eq:mpdo_bnt_identity_oblique_maps} \\
        R_\gamma
         & =V_\gamma Z_\gamma.
        \notag
    \end{align}
    Then, for every vertical letter $v$,
    \begin{align}
        L_\gamma^\dagger\widetilde{M^{(2)}}^vR_\gamma
         & =c_\gamma M_\gamma^v,
        \label{eq:mpdo_bnt_identity_oblique_compression} \\
        R_\gamma^\dagger\widetilde{M^{(2)}}^vL_\gamma
         & =c_\gamma
        (Z_\gamma^\dagger Z_\gamma)M_\gamma^v
        (Z_\gamma^\dagger Z_\gamma)^{-1}.
        \label{eq:mpdo_bnt_identity_oblique_opposite}
    \end{align}
    Write $\mc{R}(M_\gamma)^{rs}$ for the horizontal slice
    $v\mapsto(M_\gamma^v)_{rs}$. For open bond indices $r,s$ and blocked
    physical indices $i,j$, define
    \begin{align}
        f^\gamma_{(r,s),(i,j)}
         & =c_\gamma^{-1}
        \overline{(L_\gamma)_{ir}}(R_\gamma)_{js}.
        \label{eq:mpdo_bnt_identity_physical_coefficients}
    \end{align}
    These coefficients express every horizontal slice of $M_\gamma$ as a
    linear combination of the physical letters of the two-site blocking:
    \begin{align}
        \sum_{i,j}f^\gamma_{(r,s),(i,j)}(M^{[2]})^{ij}
         & =\mc{R}(M_\gamma)^{rs}.
        \label{eq:mpdo_bnt_identity_physical_span}
    \end{align}
    This gives the inclusion in the physical-letter span used to compare
    the two vertical forms in
    \cite[Appendix~C.4, lines~2048--2057]{Cirac2016MPDO_arXiv}.

    In general $L_\gamma\ne R_\gamma$, so
    (\ref{eq:mpdo_bnt_identity_oblique_compression}) is not an isometric
    compression onto a positive multiple of $M_\gamma$.
    At the marked-chain level, adjoint reflection exchanges
    the two maps together with the bond-pair and tail reversal. The reflected
    orientation is therefore governed by
    (\ref{eq:mpdo_bnt_identity_oblique_opposite}), rather than by a second copy
    of $M_\gamma$. The linear-span identity alone gives neither a common
    reflected tensor nor the Gram identity
    \begin{align}
        (Z_\gamma^\dagger Z_\gamma)M_\gamma^v
        (Z_\gamma^\dagger Z_\gamma)^{-1}
         & =M_\gamma^v.
        \notag
    \end{align}
    This prevents the direct two-site comparison from proving the Gram
    identity. The theorem below instead compares the one-site compression
    onto $M_\gamma$ with the two-site compression onto
    $Z_\gamma M_\gamma Z_\gamma^{-1}$, keeping the same unblocked tail.
\end{theorem}

\begin{proof}\leanok
    Taking adjoints in the left factor of
    (\ref{eq:mpdo_bnt_identity_oblique_maps}) and substituting
    (\ref{eq:mpdo_bnt_exact_sector_gauge_blocked_corner}) gives
    \begin{align}
        L_\gamma^\dagger\widetilde{M^{(2)}}^vR_\gamma
         & =Z_\gamma^{-1}
        (V_\gamma^\dagger\widetilde{M^{(2)}}^vV_\gamma)
        Z_\gamma
        =c_\gamma M_\gamma^v.
        \notag
    \end{align}
    Reversing the two maps instead gives
    (\ref{eq:mpdo_bnt_identity_oblique_opposite}). Finally, taking the
    $(r,s)$ matrix entry of a general two-sided compression yields
    \begin{align}
        [L^\dagger\widetilde{M^{(2)}}^vR]_{rs}
         & =\sum_{i,j}\overline{L_{ir}}R_{js}(M^{[2]})^{ij}_v.
        \notag
    \end{align}
    Divide by $c_\gamma$ and use
    (\ref{eq:mpdo_bnt_identity_oblique_compression}) to obtain
    (\ref{eq:mpdo_bnt_identity_physical_span}).
\end{proof}

\begin{theorem}[Normalization of conjugation by a scalar Gram matrix]
    \label{thm:normalized_conjugation_of_scalar_gram}
    \lean{Matrix.normalized_conj_eq_conj_inv_of_gram_eq_smul_one}
    \leanok
    If an invertible matrix $Z$ satisfies
    $Z^\dagger Z=\omega\Id$ with $\omega>0$, then conjugation by
    $\omega^{-1/2}Z$ equals conjugation by $Z$.
\end{theorem}

\begin{proof}\leanok
    The Gram identity identifies $Z^{-1}$ with
    $\omega^{-1}Z^\dagger$; hence, for every matrix $B$ of the same size,
    $(\omega^{-1/2}Z)B(\omega^{-1/2}Z)^\dagger
        =ZBZ^{-1}$.
\end{proof}

\begin{theorem}[Unitary sector conjugacy from a positive scalar Gram matrix]
    \label{thm:mpdo_bnt_conditional_identity_marked_unitary_conjugacy}
    \lean{MPOTensor.BNTAlgebraTensorClause.TwoSiteExactSectorGauge.%
        exists_unitary_sector_conjugacy_of_gauge_gram_eq_pos_smul_one}
    \leanok
    \uses{def:mpdo_bnt_two_site_exact_sector_gauge}
    Let $Z_\gamma$ be an exact invertible sector gauge and suppose
    $Z_\gamma^\dagger Z_\gamma=\omega_\gamma\Id$ with
    $\omega_\gamma>0$. Define $U_\gamma=\omega_\gamma^{-1/2}Z_\gamma$.
    Then $U_\gamma$ is unitary and, for every vertical letter $v$,
    \begin{align}
        A_{\sigma(\gamma)}^v
         & =U_\gamma M_\gamma^v U_\gamma^\dagger.
        \label{eq:mpdo_bnt_conditional_unitary_conjugacy}
    \end{align}
    This is the inverse-square-root normalization used in
    \cite[Proposition~4.13, lines~1903--1908]{Cirac2016MPDO_arXiv}, applied
    at Appendix~C.4, lines~2053--2057.
\end{theorem}

\begin{proof}\leanok
    \uses{thm:normalize_scalar_gram_matrix, thm:normalized_conjugation_of_scalar_gram}
    The positive scalar Gram identity gives
    $U_\gamma^\dagger U_\gamma=\Id$ and
    $Z_\gamma^{-1}=\omega_\gamma^{-1}Z_\gamma^\dagger$.
    Consequently
    $U_\gamma M_\gamma^vU_\gamma^\dagger
        =Z_\gamma M_\gamma^vZ_\gamma^{-1}$.
    The exact sector conjugacy then gives
    (\ref{eq:mpdo_bnt_conditional_unitary_conjugacy}).
\end{proof}

\begin{theorem}[Gram identity and unitary sector conjugacy]
    \label{thm:mpdo_bnt_mixed_prefix_unitary_sector_conjugacy}
    \lean{MPOTensor.BNTAlgebraTensorClause.TwoSiteExactSectorGauge.%
        gramDressing_gauge_eq_one}
    \lean{MPOTensor.BNTAlgebraTensorClause.TwoSiteExactSectorGauge.%
        gauge_gram_eq_pos_smul_one}
    \lean{MPOTensor.BNTAlgebraTensorClause.TwoSiteExactSectorGauge.%
        exists_unitary_sector_conjugacy}
    \lean{MPOTensor.BNTAlgebraTensorClause.TwoSiteExactSectorGauge.%
        UnitarySectorConjugacy.ofAlgebraTensorClause}
    \leanok
    \uses{def:mpdo, def:cpsv_canonical_form, def:mpdo_bnt_two_site_exact_sector_gauge}
    Let $M$ generate MPDOs, and suppose that its MPS tensor obtained by
    pairing the physical indices is in canonical form. Let the one-site and two-site
    sectors be paired by the exact gauges $Z_\gamma$ above.
    Then, for every sector $\gamma$ and every vertical letter $v$,
    \begin{align}
        (Z_\gamma^\dagger Z_\gamma)M_\gamma^v
        (Z_\gamma^\dagger Z_\gamma)^{-1}
         & =M_\gamma^v.
        \label{eq:mpdo_bnt_mixed_prefix_gram_identity}
    \end{align}
    Consequently there is a real number $\omega_\gamma>0$ such that
    \begin{align}
        Z_\gamma^\dagger Z_\gamma
         & =\omega_\gamma\Id,
        \label{eq:mpdo_bnt_mixed_prefix_scalar_gram}
    \end{align}
    and there is a unitary $U_\gamma$ satisfying
    \begin{align}
        A_{\sigma(\gamma)}^v
         & =U_\gamma M_\gamma^vU_\gamma^\dagger.
        \label{eq:mpdo_bnt_mixed_prefix_unitary_conjugacy}
    \end{align}

    The proof compares a one-site compression onto $M_\gamma$ with a
    two-site compression onto $Z_\gamma M_\gamma Z_\gamma^{-1}$,
    keeping the same unblocked tail. This supplies the common comparison
    missing from the direct application of
    \cite[Proposition~4.13, lines~1909--1919]{Cirac2016MPDO_arXiv}
    in~\cite[Appendix~C.4, lines~2048--2057]{Cirac2016MPDO_arXiv}.
    The correction to the direct two-site argument is recorded in
    \cite{gap:cpsv16_two_site_sector_unitary_gauge_gap}.
\end{theorem}

\begin{proof}\leanok
    \uses{thm:mpdo_bnt_exact_sector_one_site_raw_corner, thm:mpdo_bnt_exact_sector_gauge_blocked_corner, thm:mpdo_bnt_exact_sector_identity_oblique_physical_span, thm:mpdo_two_site_prefix_gram_reflected_marked_chain, thm:mpdo_reflected_marked_chain, thm:cpsv_original_coordinate_marked_lemma_l, thm:normal_tensor_gram_rigidity, thm:mpdo_bnt_conditional_identity_marked_unitary_conjugacy}
    Put $G_\gamma=Z_\gamma^\dagger Z_\gamma$. The two-site initial-block
    identity applied to the compression onto
    $Z_\gamma M_\gamma Z_\gamma^{-1}$ identifies the marked chains of
    $G_\gamma M_\gamma G_\gamma^{-1}$ with the reflected-adjoint chains of
    $M_\gamma$, followed by an arbitrary unblocked tail of $M^\dagger$.
    The ordinary one-site reflected identity applied to the compression onto
    $M_\gamma$ gives the same reflected-adjoint chains.

    Both marked tensors are linear combinations of the one-site physical
    letters. Indeed, the mark $M_\gamma$ is obtained from the pair
    $(W_\gamma,W_\gamma)$, whereas the mark
    $G_\gamma M_\gamma G_\gamma^{-1}$ is obtained from the pair
    $(W_\gamma G_\gamma^\dagger,W_\gamma G_\gamma^{-1})$.
    Their marked traces agree against every positive-length tail. The
    original-coordinate form of Lemma~L therefore identifies the two marked
    tensors and proves
    (\ref{eq:mpdo_bnt_mixed_prefix_gram_identity}). Normality of
    $M_\gamma$ gives
    (\ref{eq:mpdo_bnt_mixed_prefix_scalar_gram}), and normalizing
    $Z_\gamma$ by $\omega_\gamma^{-1/2}$ gives the unitary in
    (\ref{eq:mpdo_bnt_mixed_prefix_unitary_conjugacy}).
\end{proof}
